\documentclass[10pt,letterpaper]{article}
\usepackage[utf8]{inputenc}
\usepackage[spanish,es-tabla]{babel}
\usepackage{geometry}
\geometry{margin=0.68in, top=0.85in, bottom=0.78in, headheight=36pt, headsep=14pt}
\usepackage{amsmath,amssymb}
\usepackage{xcolor}
\usepackage{graphicx}
\usepackage{tabularx}
\usepackage{booktabs}
\usepackage{colortbl}
\usepackage{array}
\usepackage{enumitem}
\usepackage{fancyhdr}
\usepackage{titlesec}
\usepackage{lastpage}
\usepackage{setspace}
\usepackage{listings}

% Tipografía institucional Icesi (Sans-Serif Helvetica)
\usepackage{helvet}
\renewcommand{\familydefault}{\sfdefault}

% Interlineado aumentado un 30% (setstretch 1.38) para mayor holgura y lectura óptima
\setstretch{1.38}

% Altura de celdas en tablas
\renewcommand{\arraystretch}{1.25}

% Paleta de Colores Oficial de los Programas Icesi
\definecolor{icesiBlue}{HTML}{5454E9}       % Azul vibrante institucional
\definecolor{icesiOrange}{HTML}{F0681A}     % Naranja acento
\definecolor{icesiGreen}{HTML}{4CB978}      % Verde acento
\definecolor{icesiPurple}{HTML}{865CF0}     % Púrpura acento
\definecolor{icesiDark}{HTML}{1A1A1A}       % Texto principal oscuro
\definecolor{icesiMuted}{HTML}{6B7280}      % Gris metadatos
\definecolor{icesiBorder}{HTML}{D1D5DB}     % Líneas y bordes
\definecolor{icesiRowLine}{HTML}{E5E7EB}    % Líneas de tabla
\definecolor{icesiLight}{HTML}{F9FAFB}      % Fondo suave
\definecolor{codeBackground}{HTML}{F3F4F6}  % Fondo código

% Configuración de Listings
\lstdefinestyle{bashstyle}{
    backgroundcolor=\color{codeBackground},
    commentstyle=\color{icesiMuted}\itshape,
    keywordstyle=\color{icesiBlue}\bfseries,
    basicstyle=\ttfamily\scriptsize,
    breaklines=true,
    frame=single,
    rulecolor=\color{icesiBorder},
    language=bash,
    aboveskip=2pt,
    belowskip=2pt
}

\lstdefinestyle{codeStyle}{
    backgroundcolor=\color{codeBackground},
    commentstyle=\color{icesiMuted}\itshape,
    keywordstyle=\color{icesiBlue}\bfseries,
    stringstyle=\color{icesiGreen},
    basicstyle=\ttfamily\scriptsize,
    breaklines=true,
    frame=single,
    rulecolor=\color{icesiBorder},
    aboveskip=2pt,
    belowskip=2pt
}

% =========================================================================
% ENCABEZADO Y PIE DE PÁGINA ESTILO PROGRAMA DE CURSO ICESI
% =========================================================================
\pagestyle{fancy}
\fancyhf{}
\fancyhead[L]{%
  \begin{minipage}[c]{2.8cm}
    \includegraphics[height=0.58cm]{logo-icesi.png}
  \end{minipage}%
  \hspace{0.15cm}%
  {\color{icesiBorder}\vrule width 0.8pt height 0.60cm}%
  \hspace{0.2cm}%
  \begin{minipage}[c]{10.5cm}
    \fontsize{7.5}{9.5}\selectfont\textbf{\color{icesiDark}UNIVERSIDAD ICESI} \\
    \fontsize{7}{9}\selectfont\textbf{\color{icesiDark}PROGRAMA DE CURSO / GUÍA PRÁCTICA FORMATIVA: INGENIERÍA DE SOFTWARE V} \\
    \fontsize{6.5}{8.5}\selectfont\color{icesiMuted}Período: 202620 $\bullet$ NRC: 11005 $\bullet$ Semana 10 $\bullet$ Duración: 2 Horas
  \end{minipage}%
}
\fancyhead[R]{\fontsize{7.5}{9.5}\selectfont\color{icesiMuted}Pág. \thepage\ de \pageref{LastPage}}

% Barra horizontal multicolor característica del programa de curso
\renewcommand{\headrule}{%
  \vspace{0.08cm}%
  \noindent%
  \textcolor{icesiBlue}{\rule{0.60\textwidth}{1.8pt}}%
  \textcolor{icesiOrange}{\rule{0.18\textwidth}{1.8pt}}%
  \textcolor{icesiGreen}{\rule{0.11\textwidth}{1.8pt}}%
  \textcolor{icesiPurple}{\rule{0.11\textwidth}{1.8pt}}%
}
\renewcommand{\footrulewidth}{0.4pt}
\fancyfoot[L]{\fontsize{7.5}{9.5}\selectfont\color{icesiMuted}Ingeniería de Software V (09776 - TIC) $\bullet$ Taller Formativo: Pruebas E2E con Playwright}
\fancyfoot[R]{\fontsize{7.5}{9.5}\selectfont\color{icesiMuted}Universidad Icesi $\bullet$ DCSI}

% =========================================================================
% COMANDOS DE SECCIONES CON BADGE ESTILO PROGRAMA ICESI
% =========================================================================
\newcommand{\sectionBadge}[1]{%
  \refstepcounter{section}%
  \vspace{0.16cm}%
  \noindent%
  \colorbox{icesiBlue}{\color{white}\rule[-0.08cm]{0pt}{0.42cm}\hspace{0.16cm}\textbf{\normalsize \thesection}\hspace{0.16cm}}%
  \hspace{0.18cm}%
  {\fontsize{9.5}{12}\selectfont\textbf{\color{icesiDark}#1}}\\[-0.10cm]%
  \noindent{\color{icesiDark}\rule{\textwidth}{0.8pt}}\vspace{0.06cm}%
}

\titleformat{\subsection}
  {\color{icesiBlue}\normalfont\normalsize\bfseries}
  {\thesubsection}{0.4em}{}
\titlespacing*{\subsection}{0pt}{0.4ex plus 0.2ex minus .1ex}{0.2ex plus .1ex}

\newcolumntype{C}[1]{>{\centering\arraybackslash}p{#1}}
\newcolumntype{L}[1]{>{\raggedright\arraybackslash}p{#1}}
\newcolumntype{Y}{>{\centering\arraybackslash}X}

\begin{document}

% =========================================================================
% TÍTULO Y PRESENTACIÓN INICIAL (PÁGINA 1)
% =========================================================================
\noindent
{\fontsize{8.5}{11}\selectfont\textbf{\color{icesiBlue}GUÍA PRÁCTICA FORMATIVA: PRUEBAS END-TO-END (E2E) CON PLAYWRIGHT}}\\[0.06cm]
{\fontsize{15}{18}\selectfont\textbf{Ingeniería de Software V --- Semana 10 (Período 2026-2)}}\\[0.10cm]
{\fontsize{8.5}{11}\selectfont\color{icesiMuted}\textbf{Cúspide de la Pirámide: Auto-Waiting, Page Object Model (POM), Network Mocking y CI Headless}}\\[0.04cm]
{\color{icesiRowLine}\hrule height 0.6pt}\vspace{0.10cm}

\noindent
{\color{icesiBlue}\textbf{\rule{2.5pt}{8.5pt}}}~\textbf{\fontsize{8.5}{11}\selectfont FACULTAD DE INGENIERÍA, DISEÑO Y CIENCIAS APLICADAS} \hfill
{\color{icesiOrange}\textbf{\rule{2.5pt}{8.5pt}}}~\textbf{\fontsize{8.5}{11}\selectfont DEPARTAMENTO DE COMPUTACIÓN Y SISTEMAS INTELIGENTES}\\[0.06cm]
{\color{icesiRowLine}\hrule height 0.6pt}\vspace{0.10cm}

% =========================================================================
% TABLA DE INFORMACIÓN Y METADATOS FORMATIVOS
% =========================================================================
\noindent
{\color{icesiBlue}\textbf{\rule{2.5pt}{8.5pt}}}~\textbf{\fontsize{8.5}{11}\selectfont INFORMACIÓN DE LA ACTIVIDAD} \hfill {\fontsize{7.5}{9}\selectfont\color{icesiMuted}\textbf{IDENTIFICACIÓN}}\\[-0.22cm]
{\color{icesiBorder}\hrule height 0.5pt}\vspace{0.05cm}

\noindent
\small
\begin{tabularx}{\textwidth}{>{\bfseries\color{icesiDark}\fontsize{7.5}{9.5}\selectfont}L{2.5cm}X!{\color{icesiBlue}\vrule width 0.8pt}>{\bfseries\color{icesiDark}\fontsize{7.5}{9.5}\selectfont}L{2.5cm}X}
CÓDIGO: & 09776 - TIC & PERÍODO: & 202620 (Semestre 2026-2) \\ \hline
\rowcolor{icesiLight}
NRC / GRUPO: & 11005 / Grupo 001 & PONDERACIÓN: & \textbf{Formativa} (Sin calificación sumativa) $\bullet$ Duración: 2 Horas \\ \hline
POLÍTICA IAG: & Nivel 3: Colaboración Asistida & MODALIDAD: & Parejas / Práctica de Laboratorio Guiada \\ \hline
\rowcolor{icesiLight}
ESTUDIANTE 1: & \rule{4.5cm}{0.4pt} & CÓDIGO: & \rule{2.2cm}{0.4pt} \\ \hline
ESTUDIANTE 2: & \rule{4.5cm}{0.4pt} & CÓDIGO: & \rule{2.2cm}{0.4pt} \\ \hline
\end{tabularx}

\vspace{0.10cm}

% =========================================================================
% SECCIÓN 1: CONTEXTO Y DESCRIPCIÓN DEL RETO
% =========================================================================
\sectionBadge{Contexto del Problema \& Aseguramiento de Calidad E2E}

\subsection{Descripción del Reto de Ingeniería de Software}
En las sesiones anteriores consolidamos la base de la pirámide de calidad: pruebas unitarias con JUnit 5, cobertura con JaCoCo y análisis estático SAST con SonarQube dentro de pipelines de Jenkins. No obstante, una aplicación puede tener 100\% de cobertura unitaria y fallar ante el usuario real debido a incompatibilidades de navegador, renderizado CSS asíncrono o ruptura en los flujos de interacción del cliente web.

En esta guía práctica formativa, los estudiantes implementarán una suite profesional de pruebas \textbf{End-to-End (E2E)} con \textbf{Playwright} sobre la aplicación \textbf{Product Management CRUD} (React 18 + TypeScript + API REST). El objetivo es aprender a estructurar pruebas resilientes aplicando:
\begin{enumerate}[itemsep=1pt]
  \item \textbf{Auto-waiting y locators de accesibilidad (W3C/ARIA):} Eliminar pausas fijas arbitrarias (\texttt{sleep}) usando selectores semánticos orientados al usuario.
  \item \textbf{Page Object Model (POM):} Aislar la lógica de interacción con la UI de las aserciones de prueba para garantizar mantenibilidad.
  \item \textbf{Network Mocking:} Interceptar llamadas HTTP (\texttt{page.route}) para validar la respuesta de la interfaz ante caídas del servidor sin alterar el backend.
  \item \textbf{Validación opcional en CI/CD:} Ejecución desatendida (\textit{headless}) en GitHub Actions con publicación de trazas forenses en caso de fallo.
\end{enumerate}

\vspace{0.08cm}
\noindent
\small
\begin{tabularx}{\textwidth}{|L{2.8cm}|X|L{3.5cm}|}
\hline
\rowcolor{icesiLight}
\textbf{Componente} & \textbf{Tecnología \& Stack} & \textbf{Puerto / Endpoint} \\ \hline
\textbf{Aplicación Web (SPA)} & React 18, Vite, TypeScript, TailwindCSS & Puerto \texttt{3000} (\texttt{http://localhost:3000}) \\ \hline
\textbf{API REST Backend} & Node Mock Server / Spring Boot REST & Puerto \texttt{8080} (\texttt{/api/v1/products}) \\ \hline
\textbf{Suite E2E Playwright} & Playwright Test Runner, TypeScript & Motores: Chromium, Firefox, WebKit \\ \hline
\textbf{CI/CD (Opcional)} & GitHub Actions (\texttt{e2e-playwright.yml}) & Modo Headless + HTML Report Artifacts \\ \hline
\end{tabularx}

\newpage

% =========================================================================
% PÁGINA 2: PROCEDIMIENTO EXPERIMENTAL PASO A PASO
% =========================================================================
\sectionBadge{Procedimiento Experimental Paso a Paso}

\subsection{Fase 1: Puesta en Marcha del Código Base y Tooling (15 min)}
Ubíquese en el directorio del código base suministrado (\texttt{codigo\_base/}) e instale las dependencias de la aplicación y de Playwright:
\begin{lstlisting}[style=bashstyle]
cd codigo_base
# 1. Instalar dependencias del proyecto y Playwright
npm install

# 2. Descargar los navegadores oficiales de Playwright (Chromium, Firefox, WebKit)
npx playwright install --with-deps

# 3. Iniciar la aplicacion web y el servidor mock en segundo plano:
npm run dev
\end{lstlisting}
Verifique en su navegador local que la aplicación responde correctamente en \texttt{http://localhost:3000}.

\subsection{Fase 2: Configuración Multi-Navegador en \texttt{playwright.config.ts} (15 min)}
Inspeccione el archivo de configuración central \texttt{playwright.config.ts} asegurando los siguientes parámetros de ejecución:
\begin{lstlisting}[style=codeStyle]
import { defineConfig, devices } from '@playwright/test';

export default defineConfig({
  testDir: './e2e/tests',
  fullyParallel: true,
  retries: 1,                          // 1 reintento para aislar flakiness de red
  reporter: [['html', { open: 'never' }], ['list']],
  use: {
    baseURL: 'http://localhost:3000',
    trace: 'on-first-retry',           // Capturar traza forense al fallar
    screenshot: 'only-on-failure',
  },
  projects: [
    { name: 'chromium', use: { ...devices['Desktop Chrome'] } },
    { name: 'firefox', use: { ...devices['Desktop Firefox'] } },
    { name: 'webkit', use: { ...devices['Desktop Safari'] } },
  ],
});
\end{lstlisting}

\subsection{Fase 3: Implementación del Page Object Model (POM) (30 min)}
En \texttt{e2e/pages/ProductPage.ts}, construya la clase que encapsula los locators y métodos de interacción con el catálogo de productos. Utilice selectores de accesibilidad semántica:
\begin{lstlisting}[style=codeStyle]
import { Page, Locator } from '@playwright/test';

export class ProductPage {
  readonly page: Page;
  readonly nameInput: Locator;
  readonly priceInput: Locator;
  readonly submitButton: Locator;
  readonly successAlert: Locator;

  constructor(page: Page) {
    this.page = page;
    this.nameInput = page.getByLabel('Nombre del Producto');
    this.priceInput = page.getByLabel('Precio');
    this.submitButton = page.getByRole('button', { name: 'Registrar Producto' });
    this.successAlert = page.getByRole('alert');
  }

  async goto() {
    await this.page.goto('/');
  }

  async createProduct(name: string, price: string) {
    await this.nameInput.fill(name);
    await this.priceInput.fill(price);
    await this.submitButton.click();
  }
}
\end{lstlisting}

\subsection{Fase 4: Pruebas de Interacción y Network Mocking (30 min)}
En \texttt{e2e/tests/products.spec.ts}, escriba las pruebas para el flujo feliz y para la simulación de caída HTTP 500:
\begin{lstlisting}[style=codeStyle]
import { test, expect } from '@playwright/test';
import { ProductPage } from '../pages/ProductPage';

test.describe('Gestion de Productos - E2E Suite', () => {
  test('Flujo Feliz: Crear producto y verificar en la tabla', async ({ page }) => {
    const productPage = new ProductPage(page);
    await productPage.goto();

    await productPage.createProduct('Monitor 4K', '1250000');

    // Aserciones Web-First con auto-retry
    await expect(productPage.successAlert).toHaveText('Producto creado exitosamente');
    await expect(page.getByRole('cell', { name: 'Monitor 4K' })).toBeVisible();
  });

  test('Resiliencia: Manejar error 500 del backend mediante Network Mocking', async ({ page }) => {
    // Interceptar la llamada POST a la API REST
    await page.route('**/api/v1/products', async route => {
      if (route.request().method() === 'POST') {
        await route.fulfill({
          status: 500,
          contentType: 'application/json',
          body: JSON.stringify({ message: 'Error interno en la base de datos' })
        });
      } else {
        await route.continue();
      }
    });

    const productPage = new ProductPage(page);
    await productPage.goto();
    await productPage.createProduct('Teclado Mecanico', '350000');

    // Validar que el frontend muestra alerta de error y no se congela
    await expect(page.getByRole('alert')).toContainText('Error al procesar el producto');
  });
});
\end{lstlisting}

\subsection{Fase 5: Diagnóstico con Trace Viewer y Ejecución Headless (15 min)}
Ejecute la suite completa y examine las herramientas de diagnóstico:
\begin{lstlisting}[style=bashstyle]
# Ejecutar todas las pruebas en modo headless
npx playwright test

# En caso de fallo o inspeccion visual, abrir el reporte HTML interactivo:
npx playwright show-report
\end{lstlisting}

\subsection{Fase 6 (OPCIONAL): Validación Automatizada en Pipeline CI/CD}
Para los equipos que deseen integrar esta verificación en su pipeline de integración continua, cree el archivo \texttt{.github/workflows/e2e-playwright.yml}:
\begin{lstlisting}[style=codeStyle]
name: E2E Playwright Suite
on: [push, pull_request]

jobs:
  test-e2e:
    runs-on: ubuntu-latest
    steps:
      - uses: actions/checkout@v4
      - uses: actions/setup-node@v4
        with: { node-version: '20' }
      - name: Install dependencies & Playwright Browsers
        run: |
          npm ci
          npx playwright install --with-deps
      - name: Run E2E Tests Headless
        run: npm run test:e2e
      - name: Upload HTML Report on Failure
        uses: actions/upload-artifact@v4
        if: failure()
        with:
          name: playwright-report
          path: playwright-report/
\end{lstlisting}

% =========================================================================
% PÁGINA 3: CUESTIONARIO Y CHECKLIST FORMATIVO
% =========================================================================
\sectionBadge{Cuestionario de Reflexión Conceptual \& Checklist Formativo}

\subsection{Preguntas de Reflexión Técnica}
\begin{enumerate}[leftmargin=*,itemsep=3pt]
  \item \textbf{Actionability y Auto-Waiting:} ¿De qué manera los mecanismos internos de Playwright (comprobación de visibilidad, estabilidad y eventos) eliminan la necesidad de usar \texttt{sleep()} arbitrarios en las pruebas?
  \item \textbf{Locators W3C/ARIA vs XPath:} ¿Por qué escribir selectores basados en accesibilidad (\texttt{getByRole}, \texttt{getByLabel}) no solo hace las pruebas más resilientes a cambios de diseño, sino que promueve mejores prácticas de accesibilidad web?
  \item \textbf{Network Mocking vs Pruebas de Integración Reales:} ¿En qué casos de prueba es conveniente interceptar peticiones de red y cuándo es obligatorio probar el backend real conectado a la base de datos?
\end{enumerate}

\subsection{Checklist de Logro Formativo (Autoevaluación)}
\small
\begin{tabularx}{\textwidth}{|L{4.2cm}|X|C{2.5cm}|}
\hline
\rowcolor{icesiLight}
\textbf{Objetivo Formativo} & \textbf{Evidencia Observable} & \textbf{Estado} \\ \hline
\textbf{1. Configuración Multi-Motor} & Playwright ejecuta pruebas limpiamente en Chromium, Firefox y WebKit. & [ ] Logrado \\ \hline
\textbf{2. Page Object Model (POM)} & Selectores y acciones encapsulados en clase \texttt{ProductPage.ts}. & [ ] Logrado \\ \hline
\textbf{3. Network Mocking} & Intercepción de fallo HTTP 500 y verificación de respuesta en la UI. & [ ] Logrado \\ \hline
\textbf{4. Diagnóstico Forense} & Inspección de reporte HTML y trazas paso a paso con Trace Viewer. & [ ] Logrado \\ \hline
\textbf{5. CI/CD (Opcional)} & Pipeline en GitHub Actions configurado con reporte en artefactos. & [ ] Opcional \\ \hline
\end{tabularx}

\end{document}
